\section{The gauge-algebra restriction at every rank}
\label{sec:lie-classification}

The restriction on the simple gauge algebra is uniform in its rank.  The
representation-theoretic input below is classical; equation~(18) of
\cite{UrichukWalton} records its more general highest-weight version.  We
give the argument needed here and then determine the entire first-order
cohomology in the absence of a symmetric cubic invariant.

\begin{proposition}[Invariant cubic tensors]
\label{prop:uniform-cubic}
For a complex simple Lie algebra,
\begin{equation}
\dim(\g^{*\otimes3})^{\g}=
\begin{cases}
2,&\g=\mathfrak{sl}_N,\quad N\ge3,\\
1,&\text{every other simple }\g.
\end{cases}
\label{eq:uniform-cubic}
\end{equation}
In every case the alternating part is spanned by
$f(X,Y,Z)=\kappa(X,[Y,Z])$.  For $\mathfrak{sl}_N$, $N\ge3$, the additional
tensor is the symmetric tensor $\tr(X\{Y,Z\})$.
\end{proposition}

\begin{proof}
Let $\theta$ be the highest root, $\rho$ the half-sum of positive roots,
and $W$ the Weyl group.  Write $m_{\g}$ for the weight multiplicity in
the adjoint representation.  The invariant form identifies the adjoint
representation with its dual.  The Weyl character formula therefore gives
\begin{align}
\dim(\g^{*\otimes3})^{\g}
&=\dim\operatorname{Hom}_{\g}(\g,\g\otimes\g)\notag\\
&=\sum_{w\in W}\det(w)\,
m_{\g}\bigl(\theta+\rho-w(\theta+\rho)\bigr).
\label{eq:adjoint-weyl-sum}
\end{align}
Here $m_{\g}(0)=\operatorname{rank}\g$, every root has multiplicity one,
and all other weights have multiplicity zero.  Since $\theta+\rho$ lies
strictly inside the dominant chamber, only $w=1$ gives the zero weight.

Suppose that $\alpha=\theta+\rho-w(\theta+\rho)$ is a root.
Dominance implies that $\alpha$ is positive.  Preservation of the inner
product by $W$ implies
\begin{equation}
2(\theta+\rho,\alpha)=(\alpha,\alpha),\qquad
\langle\theta+\rho,\alpha^\vee\rangle=1.
\end{equation}
The positive coroot $\alpha^\vee$ has positive integral coefficients in
the basis of simple coroots.  Each such coefficient contributes at least
one to its pairing with $\rho$, and its pairing with the dominant
weight $\theta$ is nonnegative.  Consequently $\alpha$ is a simple root
with $\langle\theta,\alpha^\vee\rangle=0$.  Conversely each such simple
root contributes, with $w=s_\alpha$ and sign $-1$.  Regularity of
$\theta+\rho$ ensures that no second Weyl element gives the same image.
Thus
\begin{equation}
\dim(\g^{*\otimes3})^{\g}
=\operatorname{rank}\g-
\#\{i:\langle\theta,\alpha_i^\vee\rangle=0\}
=\#\{i:\langle\theta,\alpha_i^\vee\rangle\ne0\}.
\label{eq:highest-root-support}
\end{equation}
In type $A_r$ with $r\ge2$ the highest root is the sum of the first and
last fundamental weights.  In type $A_1$ and type $C_r$ it is twice one
fundamental weight.  In types $B_r$ with $r\ge3$, $D_r$ with $r\ge4$,
and all exceptional types it is one fundamental weight; $B_2=C_2$ is
already covered.  This proves~\eqref{eq:uniform-cubic}.
The tensor $f$ is nonzero and alternating.  In type $A_r$, $r\ge2$,
$\tr(X\{Y,Z\})$ is nonzero and symmetric, so these two tensors exhaust
the invariant space.  Outside this family $f$ exhausts it.
\end{proof}

\begin{lemma}[Polynomial cocycles in the alternating color channel]
\label{lem:f-channel-all-degrees}
Suppose $\g$ is simple and $(\g^{*\otimes3})^{\g}=\mathbb C f$.
In the conserved-momentum complex $\Cg$,
\begin{equation}
H^3(\Cg)_{(0,0)}=\mathbb C[s_0],\qquad
H^3(\Cg)_{(a,b)}=0\quad\text{for }a+b>0.
\label{eq:f-channel-H3}
\end{equation}
The second statement holds at every polynomial derivative order, without
rotation or dilation equivariance.
\end{lemma}

\begin{proof}
At fixed momentum degree the cochain spaces are finite-dimensional
representations of the adjoint group.  The action of a constant
$X\in\g$ is $d\iota_X+\iota_Xd$ and hence is zero on cohomology.
Complete reducibility and exactness of the invariants functor imply that
every class has an invariant representative and every invariant boundary
has an invariant primitive.  It therefore suffices to work with invariant
cochains.

Write an invariant three-cochain as
\begin{equation}
P(p,q)\,\kappa(X,[Y,Z]),\qquad p+q+r=0.
\end{equation}
Alternation of the complete cochain is equivalent to
\begin{equation}
P(p,q)=P(q,p)=P(p,-p-q).
\label{eq:P-permutations}
\end{equation}
The three four-color contractions
\[
\begin{gathered}
\kappa([X_1,X_2],[X_3,X_4]),\qquad
\kappa([X_1,X_3],[X_2,X_4]),\\
\kappa([X_1,X_4],[X_2,X_3])
\end{gathered}
\]
have only the Jacobi relation.
For example, in any embedded $\mathfrak{sl}_2$ choose two elements with
$\kappa([X,Y],[X,Y])\ne0$.  The assignments $(X,X,Y,Y)$ and
$(X,Y,Y,X)$ make the first and third contractions independent.
The Chevalley--Eilenberg differential consequently vanishes exactly when
\begin{equation}
P(p,q)+P(p+q,r)=P(q,r)+P(p,q+r).
\label{eq:P-cocycle}
\end{equation}
The momenta $p,q,r$ in this identity are independent; the fourth momentum
is $-p-q-r$.

Set $q=0$ in~\eqref{eq:P-cocycle}.  Then $P(p,0)=P(0,0)$.
For a tangent vector $u$, let
$D_2P(v,0)[u]$ denote the directional derivative in the second argument.
Differentiating~\eqref{eq:P-cocycle} in $r$ at zero gives
\begin{equation}
D_2P(p,q)[u]=D_2P(p+q,0)[u]-D_2P(q,0)[u].
\end{equation}
Symmetry in $p,q$ shows that the polynomial one-form
$D_2P(v,0)$ is closed.  Its polynomial primitive, vanishing at zero, is
\begin{equation}
H(v)=\int_0^1D_2P(tv,0)[v]\,dt.
\end{equation}
The preceding derivative identity and evaluation at $q=0$ imply
\begin{equation}
P(p,q)=P(0,0)+H(p+q)-H(p)-H(q).
\label{eq:P-primitive}
\end{equation}
The second symmetry in~\eqref{eq:P-permutations} now gives
$H(v)+H(-v)=H(0)+H(0)=0$.  Hence the two-cochain
\begin{equation}
-H(p)\,\kappa(X,Y),\qquad p+q=0,
\end{equation}
is alternating.  With the differential convention used throughout,
its differential is
\begin{equation}
\bigl[H(p+q)-H(p)-H(q)\bigr]\kappa(X,[Y,Z]).
\end{equation}
It removes every positive-degree term of $P$.  The remaining constant
multiple of $s_0$ is not exact: a momentum-zero invariant two-tensor is
symmetric, and therefore there is no alternating invariant two-cochain
of that degree.  All constructions preserve the momentum bidegree.
\end{proof}

\begin{theorem}[Uniform two-loop obstruction]
\label{thm:uniform-no-go}
Let $\g$ be a complex simple Lie algebra outside the family
$\mathfrak{sl}_N$, $N\ge2$, with any nondegenerate invariant bilinear
form.  In $\Vht$ there is no solution of the quantum master equation
through order $\hbar^2$, for any choice of closed first-order local
counterterm.  Neither adjoint invariance of that counterterm nor
$U(2)$-equivariance is required.
\end{theorem}

\begin{proof}
By Proposition~\ref{prop:uniform-cubic} and
Lemma~\ref{lem:f-channel-all-degrees}, every closed first-order Taylor
cochain $j$ is cohomologous to a constant multiple of $s_0$.
The bracket induces a bracket on cohomology, and
$\{s_0,s_0\}=0$.  Therefore
\begin{equation}
\left[\frac12\{j,j\}\right]=0.
\end{equation}
The order-two master equation requires this class to cancel $[o_2]$.
The two-loop detector is nonzero for every simple algebra other than
$\mathfrak{sl}_2$: with Killing normalization its value on $F_{223}$ is
$2\dim\g\,(8n_4-3\dim\g)$, with the nonzero expressions displayed in
the two-loop computation.  Rescaling the invariant form multiplies the
obstruction by a nonzero scalar.  Hence $[o_2]\ne0$, a contradiction.
This uses the exact two-loop calculation, not a finite-rank extrapolation
of the first-order cohomology.
\end{proof}

\subsection{Semisimple direct sums}

Let $\g=\bigoplus_i\g_i$ be semisimple, with each $\g_i$ simple.
Invariance and $[\g_i,\g_i]=\g_i$ give
\begin{equation}
\kappa(\g_i,\g_j)=0\quad(i\ne j),\qquad
(\g^{*\otimes3})^{\g}
=\bigoplus_i(\g_i^{*\otimes3})^{\g_i}.
\label{eq:semisimple-splitting}
\end{equation}
For the second identity, every distribution of three tensor factors
over more than one ideal contains an ideal occurring exactly once.
Its contribution vanishes because $(\g_i^*)^{\g_i}=0$.

\begin{corollary}[Semisimple classification in the specified sector]
\label{cor:semisimple-classification}
Under the local heat-kernel hypotheses of
Theorem~\ref{thm:all-orders-traceless}, a semisimple theory admits an
all-order $U(2)$-equivariant quantization with counterterms in $\Vht$
if and only if
\begin{equation}
\g=\bigoplus_i\mathfrak{sl}_{N_i},\qquad N_i\ge2.
\end{equation}
If a non-type-$A$ simple ideal occurs, the obstruction already appears
at order $\hbar^2$, even without equivariance of the counterterms.
\end{corollary}

\begin{proof}
Choose an invariant representative of the first-order cohomology class.
Its Taylor cochain splits over the ideals by
\eqref{eq:semisimple-splitting}.  The inverse invariant form is block
diagonal, so contractions between cubic cochains supported on distinct
ideals vanish.  The bare connected two-loop graphs likewise have all
their color labels in one ideal.  Evaluating the order-two equation on
four arguments from any non-type-$A$ ideal therefore gives precisely
the impossible simple-algebra equation of
Theorem~\ref{thm:uniform-no-go}.  Evaluating a Chevalley--Eilenberg
coboundary on that ideal remains a coboundary.  This argument uses
splitting of invariant cubic tensors; it does not assume that arbitrary
restrictions commute with the Batalin--Vilkovisky bracket.

Conversely, take the quantization of each $\mathfrak{sl}_{N_i}$ supplied
by Theorem~\ref{thm:all-orders-traceless} and sum the actions.  Cross
brackets and cross propagator contractions vanish, so the full master
equation is the sum of the individual master equations.  If the invariant
form on a factor is $c\tr$ with $c\ne0$, replace its trace-normalized
quantum action $S(\hbar)$ by $cS(\hbar/c)$.  A connected graph contributes
$c$ from each vertex and $c^{-1}$ from each internal contraction.
Together with the counterterm orders, its coefficient of $\hbar^n$
is multiplied by $c^{1-n}$, so this replacement preserves every
coefficient of the master equation.  It handles every nondegenerate
invariant form on the semisimple algebra.
\end{proof}

\begin{remark}
The exclusion of an abelian center is essential to the necessity proof.
For a reductive algebra with center, a mixed tensor consisting of the
invariant bilinear form on a simple ideal and a linear functional on the
center is allowed.  Two such cubic vertices can contract through the
center and contribute to the four-simple-argument equation.  Thus
\eqref{eq:semisimple-splitting} does not extend to reductive algebras.
\end{remark}
